\documentclass[aps,preprint]{revtex4}%
\usepackage{amsfonts}
\usepackage{amsmath}
\usepackage{amssymb}
\usepackage{graphicx}%
\setcounter{MaxMatrixCols}{30}
%TCIDATA{OutputFilter=latex2.dll}
%TCIDATA{Version=4.10.0.2363}
%TCIDATA{CSTFile=revtex4.cst}
%TCIDATA{Created=Friday, May 20, 2005 13:40:53}
%TCIDATA{LastRevised=Saturday, May 21, 2005 17:35:02}
%TCIDATA{<META NAME="GraphicsSave" CONTENT="32">}
%TCIDATA{<META NAME="DocumentShell" CONTENT="Articles\SW\REVTeX 4">}
\newtheorem{theorem}{Theorem}
\newtheorem{acknowledgement}[theorem]{Acknowledgement}
\newtheorem{algorithm}[theorem]{Algorithm}
\newtheorem{axiom}[theorem]{Axiom}
\newtheorem{claim}[theorem]{Claim}
\newtheorem{conclusion}[theorem]{Conclusion}
\newtheorem{condition}[theorem]{Condition}
\newtheorem{conjecture}[theorem]{Conjecture}
\newtheorem{corollary}[theorem]{Corollary}
\newtheorem{criterion}[theorem]{Criterion}
\newtheorem{definition}[theorem]{Definition}
\newtheorem{example}[theorem]{Example}
\newtheorem{exercise}[theorem]{Exercise}
\newtheorem{lemma}[theorem]{Lemma}
\newtheorem{notation}[theorem]{Notation}
\newtheorem{problem}[theorem]{Problem}
\newtheorem{proposition}[theorem]{Proposition}
\newtheorem{remark}[theorem]{Remark}
\newtheorem{solution}[theorem]{Solution}
\newtheorem{summary}[theorem]{Summary}
\newenvironment{proof}[1][Proof]{\noindent\textbf{#1.} }{\ \rule{0.5em}{0.5em}}
\begin{document}
\preprint{HEP/123-qed}
\title[Short title for running header]{Long title that appears on the title page}
\author{First Author}
\affiliation{My Institution}
\author{Second Author}
\affiliation{My Institution}
\author{Third Author}
\affiliation{Other Institution}
\keywords{one two three}
\pacs{PACS number}

\begin{abstract}
Shell document for REV\TeX{} 4.

\end{abstract}
\volumeyear{year}
\volumenumber{number}
\issuenumber{number}
\eid{identifier}
\date[Date text]{date}
\received[Received text]{date}

\revised[Revised text]{date}

\accepted[Accepted text]{date}

\published[Published text]{date}

\startpage{101}
\endpage{102}
\maketitle
\tableofcontents


Consider
\begin{equation}
f\left(  x_{1},x_{2}\right)  =x_{1}^{2}-3x_{2}^{2}+6x_{1}x_{2}+3x_{1}+2x_{2}-9
\end{equation}
then
\begin{equation}
f\left(  1,1\right)  =0,
\end{equation}%
\begin{align}
\frac{\partial f\left(  x_{1},x_{2}\right)  }{\partial x_{1}} &
=2x_{1}+6x_{2}+3\\
\frac{\partial f\left(  1,1\right)  }{\partial x_{1}} &  =11
\end{align}
and
\begin{align}
\frac{\partial f\left(  x_{1},x_{2}\right)  }{\partial x_{2}} &
=-3x_{2}+6x_{1}+2\\
\frac{\partial f\left(  1,1\right)  }{\partial x_{2}} &  =5.
\end{align}
Expanding
\begin{align}
f\left(  x_{1}+\delta x_{1},x_{2}+\delta x_{2}\right)   &  =f\left(
x_{1},x_{2}\right)  +\left(
\begin{array}
[c]{cc}%
\frac{\partial f\left(  x_{1},x_{2}\right)  }{\partial x_{1}} & \frac{\partial
f\left(  x_{1},x_{2}\right)  }{\partial x_{2}}%
\end{array}
\right)  \left(
\begin{array}
[c]{c}%
\delta x_{1}\\
\delta x_{2}%
\end{array}
\right)  \\
&  +\frac{1}{2}\left(
\begin{array}
[c]{cc}%
\delta x_{1} & \delta x_{2}%
\end{array}
\right)  \left(
\begin{array}
[c]{cc}%
\frac{\partial^{2}f\left(  x_{1},x_{2}\right)  }{\partial x_{1}^{2}} &
\frac{\partial^{2}f\left(  x_{1},x_{2}\right)  }{\partial x_{1}\partial x_{2}%
}\\
\frac{\partial^{2}f\left(  x_{1},x_{2}\right)  }{\partial x_{2}\partial x_{1}}
& \frac{\partial^{2}f\left(  x_{1},x_{2}\right)  }{\partial x_{2}^{2}}%
\end{array}
\right)  \left(
\begin{array}
[c]{c}%
\delta x_{1}\\
\delta x_{2}%
\end{array}
\right)  ,
\end{align}
where
\begin{equation}
\left(
\begin{array}
[c]{cc}%
\frac{\partial^{2}f\left(  x_{1},x_{2}\right)  }{\partial x_{1}^{2}} &
\frac{\partial^{2}f\left(  x_{1},x_{2}\right)  }{\partial x_{1}\partial x_{2}%
}\\
\frac{\partial^{2}f\left(  x_{1},x_{2}\right)  }{\partial x_{2}\partial x_{1}}
& \frac{\partial^{2}f\left(  x_{1},x_{2}\right)  }{\partial x_{2}^{2}}%
\end{array}
\right)  =\left(
\begin{array}
[c]{cc}%
2 & 6\\
6 & -3
\end{array}
\right)  .
\end{equation}
To first order
\begin{equation}
f\left(  1+\delta x_{1},1+\delta x_{2}\right)  =\left(
\begin{array}
[c]{cc}%
\frac{\partial f\left(  1,1\right)  }{\partial x_{1}} & \frac{\partial
f\left(  1,1\right)  }{\partial x_{2}}%
\end{array}
\right)  \left(
\begin{array}
[c]{c}%
\delta x_{1}\\
\delta x_{2}%
\end{array}
\right)
\end{equation}
giving
\begin{equation}
\delta x_{1}=-\left[  \frac{\partial f\left(  1,1\right)  }{\partial x_{1}%
}\right]  ^{-1}\frac{\partial f\left(  1,1\right)  }{\partial x_{2}}\delta
x_{2}%
\end{equation}
Noting that
\begin{align}
\delta x_{1} &  =x_{1}-1\\
\delta x_{2} &  =x_{2}-1
\end{align}
we obtain
\begin{align}
x_{1} &  =-\left[  \frac{\partial f\left(  1,1\right)  }{\partial x_{1}%
}\right]  ^{-1}\frac{\partial f\left(  1,1\right)  }{\partial x_{2}}\left(
x_{2}-1\right)  +1\\
&  =-\left[  \frac{\partial f\left(  1,1\right)  }{\partial x_{1}}\right]
^{-1}\frac{\partial f\left(  1,1\right)  }{\partial x_{2}}x_{2}+\left[
\frac{\partial f\left(  1,1\right)  }{\partial x_{1}}\right]  ^{-1}%
\frac{\partial f\left(  1,1\right)  }{\partial x_{2}}+1\\
&  =-\frac{5}{11}x_{2}+\frac{16}{11}\equiv x_{1}\left(  x_{2}\right)
\end{align}%
\begin{align}
\frac{\partial f\left(  x_{1},x_{2}\right)  }{\partial x_{1}} &
=2x_{1}+6x_{2}+3\\
\frac{\partial f\left(  1,1\right)  }{\partial x_{1}} &  =11
\end{align}
and
\begin{align}
\frac{\partial f\left(  x_{1},x_{2}\right)  }{\partial x_{2}} &
=-3x_{2}+6x_{1}+2\\
\frac{\partial f\left(  1,1\right)  }{\partial x_{2}} &  =5.
\end{align}
Evaluating this function at $x_{2}=2$ one obtains
\begin{equation}
x_{1}\left(  2\right)  =\frac{6}{11}%
\end{equation}
the value of the function $f$ at this point is
\begin{align}
f\left(  \frac{6}{11},2\right)   &  =\frac{36}{121}-12+\frac{12}{11}+\frac
{18}{11}+4-9\\
&  =\frac{36}{121}-17+\frac{30}{11}\neq0.
\end{align}
%

\begin{align}
f\left(  x_{1}+\delta x_{1},x_{2}+\delta x_{2}\right)   &  =f\left(
x_{1},x_{2}\right)  +\left(
\begin{array}
[c]{cc}%
2x_{1}+6x_{2}+3 & -3x_{2}+6x_{1}+2
\end{array}
\right)  \left(
\begin{array}
[c]{c}%
\delta x_{1}\\
\delta x_{2}%
\end{array}
\right)  \\
&  +\frac{1}{2}\left(
\begin{array}
[c]{cc}%
\delta x_{1} & \delta x_{2}%
\end{array}
\right)  \left(
\begin{array}
[c]{cc}%
2 & 6\\
6 & -3
\end{array}
\right)  \left(
\begin{array}
[c]{c}%
\delta x_{1}\\
\delta x_{2}%
\end{array}
\right)  ,
\end{align}



\end{document}